\documentclass[11pt]{article}
% =========================================================
%  學生講義 共用前言（以 XeLaTeX 編譯）
%  需要套件：xeCJK, tcolorbox（其餘多在 BasicTeX 內）
% =========================================================
\usepackage[a4paper,margin=2.3cm]{geometry}
\usepackage{amsmath,amssymb}
\usepackage{xcolor}
\usepackage{graphicx}

% ---- 中文字型（macOS 系統字型）----
\usepackage{xeCJK}
\setCJKmainfont{Songti TC}      % 內文：宋體
\setCJKsansfont{Heiti TC}       % 標題/強調：黑體
\xeCJKsetup{AutoFakeBold=2.2, AutoFakeSlant=0.2}

% ---- 版面 ----
\setlength{\parindent}{0pt}
\setlength{\parskip}{0.55em}
\linespread{1.15}
\renewcommand{\baselinestretch}{1.15}

% ---- 顏色 ----
\definecolor{cBlue}{HTML}{1f6feb}
\definecolor{cGray}{HTML}{6b7280}
\definecolor{cGrayBg}{HTML}{f3f4f6}
\definecolor{cPurp}{HTML}{7a3ff2}
\definecolor{cPurpBg}{HTML}{f5f1ff}

% ---- 標註框 ----
\usepackage[most]{tcolorbox}

% 就地補充新概念（灰底）：\begin{補充框}{標題}
\newtcolorbox{補充框}[1]{
  enhanced, breakable, arc=2pt, boxrule=0.6pt,
  colback=cGrayBg, colframe=cGray,
  left=9pt,right=9pt,top=7pt,bottom=7pt,
  before upper={\textbf{\sffamily #1}\par\vspace{3pt}},
}

% 延伸 / 開放問題（紫色左邊條）：\begin{延伸框}{標題}
\newtcolorbox{延伸框}[1]{
  enhanced, breakable, arc=2pt, boxrule=0pt,
  colback=cPurpBg, colframe=cPurpBg,
  borderline west={3pt}{0pt}{cPurp},
  left=10pt,right=9pt,top=7pt,bottom=7pt,
  before upper={\textbf{\sffamily 延伸閱讀｜#1}\par\vspace{3pt}},
}

% ---- 圖：置中、底下小字說明 ----
% \fig{檔名}{寬度}{說明}
\newcommand{\fig}[3]{%
  \begin{center}
  \includegraphics[width=#2\linewidth]{#1}\\[2pt]
  {\small\color{cGray} #3}
  \end{center}}

% ---- 章節標題用黑體 ----
\newcommand{\h}[1]{\sffamily #1}


\begin{document}

\begin{center}
{\sffamily\Huge\bfseries 數1-2 多項式函數}\\[3pt]
{\sffamily\large 普通高中 數學（必修・第一冊）・學生講義}
\end{center}
\vspace{0.6em}

本章先把\textbf{多項式（polynomial）}當成一種可以加、減、乘、除的代數物件，建立它的\textbf{除法原理}（餘式定理、因式定理、綜合除法）；再把多項式看成\textbf{函數} $y=f(x)$，依序畫出一次、二次、三次函數的圖形，掌握它們的對稱性與形狀；最後反過來用圖形\textbf{解多項式不等式}。

多項式是後續多數單元的共同語言：配方法直接通往圓的標準式，函數圖形的觀念是三角函數、指數與對數函數的基礎，除法原理與因式分解則貫穿整個高中代數。

\medskip
本章主線：

\begin{center}
多項式與除法原理 $\rightarrow$ 一次與二次函數（配方法）$\rightarrow$ 三次函數的圖形特徵 $\rightarrow$ 多項式不等式
\end{center}

\section*{\sffamily 2-1\quad 多項式與除法原理}

\subsection*{\sffamily A.\ 多項式是什麼}

形如
$$f(x)=a_n x^n + a_{n-1}x^{n-1}+\cdots+a_1 x + a_0$$
的式子稱為\textbf{多項式（polynomial）}，其中 $a_n,\dots,a_0$ 是常數（係數），$n$ 是非負整數。當 $a_n\neq 0$ 時，$n$ 稱為這個多項式的\textbf{次數（degree）}，記作 $\deg f$；$a_n$ 稱為\textbf{領導係數（leading coefficient）}，$a_0$ 稱為\textbf{常數項}。

多項式的次數，是其中\textbf{次數最高的非零項}之指數。一次式為 $ax+b$（$a\neq 0$），二次式為 $ax^2+bx+c$，三次式為 $ax^3+\cdots$。非零常數（如 $5$）視為\textbf{零次}；數字 $0$ 不定義次數。

多項式之間可以\textbf{加、減、乘}，結果仍是多項式：加減是同次項的係數相加減；乘法是每一項兩兩相乘再合併，且 $\deg(fg)=\deg f+\deg g$。

\textbf{注意：}$\dfrac{1}{x}=x^{-1}$、$\sqrt{x}=x^{1/2}$ 不是多項式，因為次數必須是非負整數。「次數」看的是\textbf{最高次非零項}，不是項數；$x^5-2$ 是五次式（只有兩項）。

\subsection*{\sffamily B.\ 多項式的長除法}

兩個整數相除有「商與餘數」，多項式也一樣。把一個多項式 $f(x)$ 除以另一個多項式 $g(x)$（$g\neq 0$），可以唯一地寫成
$$\boxed{\,f(x)=g(x)\,q(x)+r(x)\,},\qquad r(x)=0\ \text{或}\ \deg r<\deg g.$$
這稱為\textbf{除法原理（division algorithm）}：$q(x)$ 是\textbf{商}、$r(x)$ 是\textbf{餘式}，且餘式的次數一定\textbf{低於除式}。做法是模仿整數長除法：每一步用被除式的最高次項去除除式的最高次項，得到商的一項，乘回去相減，降一次再做，直到餘式次數低於除式為止。

\subsection*{\sffamily C.\ 除以 $(x-a)$：綜合除法與餘式定理}

最常用的是除式為\textbf{一次式 $x-a$} 的情形。此時餘式次數低於 1，所以\textbf{餘式是一個常數}。除法原理變成
$$f(x)=(x-a)\,q(x)+r,\qquad r\ \text{是常數}.$$
把 $x=a$ 代進去，$(x-a)$ 那一塊歸零，立刻得到：

\begin{補充框}{餘式定理為什麼成立？}
\textbf{餘式定理（remainder theorem）}：多項式 $f(x)$ 除以 $x-a$ 的餘式，等於把 $a$ 直接代入，即
$$\boxed{\,r=f(a)\,}.$$
也就是說，要算餘式不必真的去除，代一個數進去即可。

這個定理不是新的規則，而是除法原理在 $x=a$ 這一點上的\textbf{直接推論}。除法原理
$$f(x)=(x-a)\,q(x)+r$$
對\textbf{每一個} $x$ 都成立（它是恆等式，不是方程式），餘式 $r$ 因為次數須低於除式 $x-a$（一次）而必為常數。既然每個 $x$ 都成立，就代入最方便的 $x=a$：此時 $(x-a)=(a-a)=0$，前面 $(x-a)q(x)$ 整塊歸零，只剩
$$f(a)=0\cdot q(a)+r=r,$$
於是 $r=f(a)$。可類比為整數除法 $17=5\times3+2$，差別在於 $x-a$ 可以\textbf{令它等於零}（取 $x=a$），一步把商那一項消掉，餘式就單獨顯現。

\textbf{注意：}餘式定理只在除式為 $x-a$ 時成立。若除式是 $x^2-1$ 之類的高次式，餘式可能不是常數（會是一次以下），就不能用「代一個數」直接得到，須用待定係數或長除法。
\end{補充框}

既然除以 $x-a$ 的餘式只是一個數，整個除法可以用\textbf{綜合除法（synthetic division）}省略寫 $x$，只搬係數：

\begin{補充框}{綜合除法怎麼做}
把 $f$ 的各次係數（缺項補 $0$）排成一列，左邊放 $a$，逐項「落下、乘、加」：先把最高次係數\textbf{落下}；將它\textbf{乘以} $a$，加到下一個係數上；重複到最後。最末一格即餘式 $f(a)$，其餘各格依次為商的係數（商的次數比 $f$ 少一次）。

\textbf{注意：}
\begin{itemize}\setlength{\itemsep}{1pt}
\item 綜合除法\textbf{只能用在除式為 $x-a$}（一次、領導係數為 1）。除以 $2x-1$ 須先寫成 $2\big(x-\tfrac12\big)$ 處理，或改用長除法。
\item 缺項要\textbf{補零}：缺的次方項係數寫 $0$，否則係數會錯位一格。
\item 表格最右一格就是 $f(a)$。綜合除法與餘式定理是一體兩面。
\end{itemize}
\end{補充框}

\subsection*{\sffamily D.\ 因式定理}

餘式定理有個立刻可用的特例。若代入剛好得 $0$，表示餘式為 $0$，也就是\textbf{整除}：

\begin{補充框}{因式定理為什麼成立？}
\textbf{因式定理（factor theorem）}：對多項式 $f(x)$，
$$\boxed{\,x-a\ \text{是}\ f(x)\ \text{的因式}\iff f(a)=0\,},$$
即 $a$ 是 $f$ 的一個\textbf{根（root）}，等價於 $(x-a)$ 是 $f$ 的因式。

它只是餘式定理再追問一句：什麼時候整除（餘式為 0）？由餘式定理 $r=f(a)$，所以
$$x-a\ \text{整除}\ f(x)\iff r=0\iff f(a)=0.$$
因此因式定理與餘式定理其實是同一件事的兩種說法。

\textbf{注意：}因式定理要求 $f(a)$ \textbf{恰為 $0$}，不是「$f(a)$ 很小」或「很接近 $0$」。
\end{補充框}

因式定理是\textbf{因式分解高次多項式}的主要工具：先用「湊根」找一個 $a$ 使 $f(a)=0$，再用綜合除法把 $(x-a)$ 除掉降次。整係數多項式的有理根，\textbf{分子是常數項的因數、分母是領導係數的因數}（有理根檢驗法），所以要試的範圍其實不大。

\section*{\sffamily 2-2\quad 一次與二次函數}

\subsection*{\sffamily A.\ 從方程式到 $f(x)$ 記法}

國中把 $y=2x+1$ 看成「$x,y$ 的關係」。高中改用另一個觀點：把它看成一台\textbf{機器}，輸入一個 $x$、輸出一個 $y$，記作 $f(x)=2x+1$。於是「當 $x=3$ 時 $y$ 是多少」就寫成求 $f(3)=7$。這個 $f(x)$ 記法的好處是：可以同時談多台機器（$f,g,h$），也便於表達代入、合成與比較。\textbf{多項式函數}就是「輸出由一個多項式算出來」的函數。

\subsection*{\sffamily B.\ 一次函數與斜率}

\textbf{一次函數}為
$$f(x)=mx+k\quad(m\neq 0),$$
圖形是一條\textbf{直線}：$m$ 是\textbf{斜率（slope）}，$k$ 是 $y$ 截距（直線與 $y$ 軸交點的 $y$ 值）。斜率 $m=\dfrac{\Delta y}{\Delta x}$，表示「$x$ 每增加 1，$y$ 變化 $m$」。

把 $y=mx+k$ 和最單純的 $y=mx$（過原點、同斜率）比較：兩者斜率相同、永遠平行，差別只在 $y=mx+k$ 是 $y=mx$ \textbf{整條往上（或往下）鉛直平移 $k$}。這個「同斜率只差一個鉛直平移」的觀念，在二次函數會再用一次。

\medskip
數線上 $A(a)$、$B(b)$ 兩點，\textbf{內分}線段 $\overline{AB}$ 成 $m:n$ 的分點 $P$，其坐標為
$$\boxed{\,P=\frac{n\,a+m\,b}{m+n}\,}.$$
靠近哪一端，那一端的權重就小；中點是 $m:n=1:1$，得 $\dfrac{a+b}{2}$。

\subsection*{\sffamily C.\ 二次函數與配方法}

\textbf{二次函數}為
$$f(x)=ax^2+bx+c\quad(a\neq 0),$$
圖形是\textbf{拋物線（parabola）}。$a>0$ 開口向上、$a<0$ 開口向下；$|a|$ 越大開口越窄。

一般式 $ax^2+bx+c$ 看不出頂點在哪。\textbf{配方法（completing the square）}是把它湊成「一個完全平方加上一個常數」，讓頂點顯現：
$$
\begin{aligned}
ax^2+bx+c
&= a\!\left(x^2+\tfrac{b}{a}x\right)+c\\
&= a\!\left(x^2+\tfrac{b}{a}x+\tfrac{b^2}{4a^2}\right)-a\cdot\tfrac{b^2}{4a^2}+c\\
&= a\!\left(x+\tfrac{b}{2a}\right)^2+\left(c-\tfrac{b^2}{4a}\right).
\end{aligned}
$$
於是任何二次函數都能寫成\textbf{頂點式（標準式）}：
$$\boxed{\,f(x)=a(x-h)^2+k\,},\qquad h=-\frac{b}{2a},\ \ k=f(h).$$
這代表把 $y=ax^2$（頂點在原點）\textbf{整個平移}到頂點 $(h,k)$。由標準式 $f(x)=a(x-h)^2+k$ 可直接讀出：\textbf{頂點} $(h,k)$、\textbf{對稱軸}為鉛直線 $x=h$；$a>0$ 時頂點是最低點、\textbf{最小值} $=k$（無最大值），$a<0$ 時頂點是最高點、\textbf{最大值} $=k$（無最小值）。

\begin{補充框}{配方法為什麼能找到頂點？}
配方法把「一個看不出最值的式子」改寫成「一個一定 $\ge 0$ 的平方，再加上一個固定的數」。盯著 $(x-h)^2$ 這一塊，它有一個確定的性質：
$$(x-h)^2\ge 0,\qquad \text{而且只有當}\ x=h\ \text{時才等於 }0.$$
於是分兩種情況：
\begin{itemize}\setlength{\itemsep}{1pt}
\item \textbf{$a>0$（開口向上）}：$a(x-h)^2\ge 0$，所以 $f(x)=a(x-h)^2+k\ge k$，等號在 $x=h$ 時成立。即 $f$ 的最小值是 $k$，發生在 $x=h$。
\item \textbf{$a<0$（開口向下）}：$a(x-h)^2\le 0$，所以 $f(x)\le k$，$f$ 的最大值是 $k$，一樣在 $x=h$。
\end{itemize}
這就是頂點 $(h,k)$ 的來歷：$h$ 是讓平方歸零的那個 $x$，$k$ 是平方歸零時剩下的值（最值）。配方法把「找最值」翻譯成「平方何時等於零」這個直接可答的問題。

對稱軸也隨之而出：$(x-h)^2$ 只看「$x$ 離 $h$ 多遠」，離 $h$ 一樣遠的兩個 $x$（如 $h+t$ 與 $h-t$）會給出相同的 $f$ 值，所以圖形對 $x=h$ 左右對稱。

代數步驟也不是硬湊：要把 $x^2+\tfrac{b}{a}x$ 補成完全平方，須湊 $\big(x+\tfrac{b}{2a}\big)^2=x^2+\tfrac{b}{a}x+\tfrac{b^2}{4a^2}$，即「一次項係數的一半，再平方」就是要補的數。補上再扣回以保持相等，整理即得 $h=-\tfrac{b}{2a}$。因此 $-\tfrac{b}{2a}$ 是「一次項係數一半」自然的結果。
\end{補充框}

\medskip
上面的最值都假設 $x$ 可取\textbf{任意實數}，所以最值一定發生在頂點。但定義域常被限制在一段\textbf{閉區間（closed interval）}$[p,q]$ 上（例如邊長不可能是負的，限定 $0\le x\le 6$）。這時不能直接把頂點的 $x$ 坐標代進去，須先看頂點 $x=h$ 有沒有落在區間裡。

\begin{補充框}{在閉區間 $[p,q]$ 上求二次函數的最大、最小值}
設 $f(x)=a(x-h)^2+k$，定義域限制在 $p\le x\le q$。先比較對稱軸 $h$ 與區間的相對位置，分兩種情況：

\textbf{一、頂點落在區間內（$p\le h\le q$）}：頂點是其中一個極值。
\begin{itemize}\setlength{\itemsep}{1pt}
\item \textbf{$a>0$（開口向上）}：\textbf{最小值} $=f(h)=k$（在頂點取得）；\textbf{最大值}在\textbf{離 $h$ 較遠}的那個端點取得。
\item \textbf{$a<0$（開口向下）}：\textbf{最大值} $=k$（在頂點取得）；\textbf{最小值}在離 $h$ 較遠的端點取得。
\end{itemize}

\textbf{二、頂點落在區間外（$h<p$ 或 $h>q$）}：函數在整段 $[p,q]$ 上\textbf{單調}（一路遞增或一路遞減），因此\textbf{最大、最小值都在兩端點 $f(p),\,f(q)$}，比大小即可，頂點根本取不到。此時 $a$ 的正負只影響「往哪邊單調」，最值仍只在端點之間挑。

一句話：\textbf{永遠先檢查對稱軸在不在區間內}；若不在，極值就只在兩端點之間挑。

\textbf{注意：}最常見的錯是「不管定義域，直接把頂點當答案」。若對稱軸落在區間外，頂點的 $y$ 值根本取不到，硬填就會錯。同一個二次函數換了區間，連「最小值在不在頂點」都可能不同，務必先定位 $h$ 與 $[p,q]$ 的關係，再決定最值在頂點還是端點。
\end{補充框}

\fig{d12-interval}{0.95}{圖 2-1：閉區間上的二次極值，以 $f(x)=(x-2)^2-3$（頂點 $(2,-3)$、開口向上）為例。左——定義域 $[0,3]$，對稱軸 $x=2$ 落在區間內，最小值在頂點 $f(2)=-3$、最大值在離軸較遠的端點 $f(0)=1$。右——定義域 $[3,5]$，對稱軸落在區間外，區段單調遞增，最小、最大值都在端點 $f(3)=-2$、$f(5)=6$，頂點取不到。粗藍線為定義域內的圖形，灰色虛線為區間外不計入的部分。}

\fig{d12-vertex}{0.78}{圖 2-2：以 $y=(x-2)^2-3$ 為例。把基準拋物線 $y=x^2$（灰色虛線）右移 2、下移 3，得到頂點 $(2,-3)$、對稱軸 $x=2$ 的拋物線（藍色實線）。}

\textbf{注意：}
\begin{itemize}\setlength{\itemsep}{1pt}
\item 配方時\textbf{領導係數要先提出來}：$2x^2-8x$ 須寫成 $2(x^2-4x)$ 再配，不能直接對 $2x^2$ 配。
\item 對稱軸是 $x=-\dfrac{b}{2a}$，注意\textbf{負號}與分母 $2a$。例如 $f(x)=2x^2-8x+5$ 的軸是 $x=-\tfrac{-8}{2\cdot2}=2$，不是 $x=4$。
\item \textbf{最小值是 $y$ 的值（$k$），不是發生的位置（$h$）}。$(x-h)^2$ 永遠 $\ge0$；可能為負的是前面的係數 $a$。
\end{itemize}

\section*{\sffamily 2-3\quad 三次函數的圖形特徵}

二次函數的圖形是\textbf{左右對稱}（軸對稱）的拋物線。三次函數不再軸對稱，但有另一種對稱——\textbf{點對稱}。

\subsection*{\sffamily A.\ 對稱性：二次軸對稱，三次點對稱}

二次函數 $y=ax^2$ 關於 $y$ 軸\textbf{鏡射對稱}，對稱軸是一條鉛直線。三次函數 $y=x^3$ 則是把圖形繞\textbf{原點旋轉 $180^\circ$} 會與自己重合，稱為關於原點\textbf{點對稱}：$f(-x)=-f(x)$，是\textbf{奇函數（odd function）}。一般的三次函數 $y=ax^3+bx^2+cx+d$ 也有這種點對稱，只是對稱中心不在原點，而是落在某一點 $(h,k)$ 上。

\begin{補充框}{三次函數的對稱中心}
任何三次函數 $f(x)=ax^3+bx^2+cx+d$ 都有\textbf{唯一的對稱中心}，其 $x$ 坐標為
$$\boxed{\,h=-\frac{b}{3a}\,},$$
對稱中心為 $(h,\,f(h))$，圖形關於這一點旋轉 $180^\circ$ 不變。

\textbf{注意：}三次函數的對稱中心 $x$ 坐標是 $-\dfrac{b}{3a}$（分母為 $3a$），與二次的對稱軸 $-\dfrac{b}{2a}$（分母 $2a$）不同，兩者不可混用。三次函數\textbf{沒有}鉛直對稱軸——把紙左右對折不會重合，須繞中心旋轉 $180^\circ$ 才重合。
\end{補充框}

\fig{d12-cubic}{0.95}{圖 2-3：左——$f(x)=x^3-3x$ 關於對稱中心 $(0,0)$ 點對稱：將點 $(-1,2)$ 與 $(1,-2)$ 連線，中點即對稱中心。右——大域趨勢由最高次項 $x^3$ 決定（兩端往相反方向延伸）；在對稱中心附近放大，曲線近似一條直線。}

\subsection*{\sffamily B.\ 大域看最高次項，局部近似一條直線}

三次函數的圖形有兩個關鍵特徵，作圖時可直接運用：

\textbf{一、大域（global）特徵由最高次項決定。} 當 $x$ 很大或很小時，$ax^3$ 遠大於其他項，因此
\begin{itemize}\setlength{\itemsep}{1pt}
\item $a>0$：左下到右上（$x\to-\infty$ 時 $y\to-\infty$，$x\to+\infty$ 時 $y\to+\infty$）。
\item $a<0$：左上到右下。
\end{itemize}
這也是為何三次函數的圖形\textbf{兩端必往相反方向延伸}，因此\textbf{至少穿過 $x$ 軸一次}——每個三次方程式至少有一個實根。

\textbf{二、局部（local）近似一條直線。} 把圖形放大看任一小段，曲線近似為直的（這是高二微分「切線」的前奏）。在對稱中心附近，這條近似直線的傾斜由一次項決定。

\medskip
形狀分兩種情況（以 $a>0$ 為例）：若有兩個轉折（一個局部極大、一個局部極小），圖形呈「上、下、上」的 $N$ 形，可能與 $x$ 軸交\textbf{三}點；若沒有轉折（單調遞增），圖形一路往上，只與 $x$ 軸交\textbf{一}點。兩種情況都關於對稱中心點對稱。

\textbf{注意：}「局部近似直線」不代表它\textbf{是}直線；只是放大某一小段看起來像直線，整體仍是彎曲的三次曲線。

\begin{延伸框}{為什麼每個三次函數都有對稱中心}
二次有對稱軸還算直覺，但「每一個三次函數都恰有一個對稱中心」並不明顯。理由是：\textbf{只要把座標原點搬到對稱中心，三次函數就會顯露它的奇函數本性。}

先看最單純的 $y=x^3$，它滿足 $f(-x)=(-x)^3=-x^3=-f(x)$，是奇函數，圖形關於原點點對稱（繞原點轉 $180^\circ$ 不變）。一般的 $f(x)=ax^3+bx^2+cx+d$ 之所以也有這種對稱，是因為它\textbf{只是 $x^3$ 型曲線平移後的樣子}，而平移不會破壞對稱，只會把對稱中心一起搬走。

具體做法：令 $x=t+h$（把新原點移到 $x=h$ 處），目標是讓\textbf{二次項消失}。展開後 $t^2$ 的係數會是 $3ah+b$，要它為零：
$$3ah+b=0 \;\Rightarrow\; h=-\frac{b}{3a},$$
這正是課本給的數值。取此 $h$ 後，函數（相對於新原點）變成
$$g(t)=a t^3 + p t + r,$$
已經\textbf{沒有 $t^2$ 項}。再把常數項 $r$ 吸收掉（把新原點高度也對齊到中心），剩下的 $at^3+pt$ 滿足 $a(-t)^3+p(-t)=-(at^3+pt)$，是\textbf{奇函數}，關於新原點點對稱。把新原點換回原座標，即得對稱中心 $\big(-\tfrac{b}{3a},\,f(-\tfrac{b}{3a})\big)$。每個三次函數都能做這個平移，所以每個都有對稱中心。

這與二次函數配方「把一次項擠成完全平方、找出對稱軸」是同一種思路：\textbf{用平移把式子化簡到能看出對稱}。二次配方後是偶函數 $at^2$（鏡射對稱），三次平移後是奇函數 $at^3+pt$（點對稱）。這不是開放問題，而是\textbf{確定的代數事實}。至於四次函數，一般\textbf{沒有}單一對稱中心。
\end{延伸框}

\section*{\sffamily 2-4\quad 多項式不等式}

學會畫圖之後，解不等式就成為「看圖」的事。$f(x)>0$ 問的是「圖形在 $x$ 軸\textbf{上方}的 $x$ 範圍」，$f(x)<0$ 問的是「圖形在 $x$ 軸\textbf{下方}的範圍」。關鍵在於：\textbf{多項式函數的圖形只有在根（與 $x$ 軸的交點）處才會換邊}（由正變負或由負變正）。所以「找根 $\rightarrow$ 把數線切成幾段 $\rightarrow$ 每段判斷正負」即可求解。

\subsection*{\sffamily A.\ 一次不等式}

一次式 $ax+b$ 只有一個根，圖形是斜直線，過根之後變號，直接移項求解即可。

\textbf{注意：}\textbf{兩邊同乘（或同除）一個負數，不等號要反向}。例如 $-2x>6$ 兩邊除以 $-2$ 得 $x<-3$（不是 $x>-3$）。這是一次不等式主要的陷阱。

\subsection*{\sffamily B.\ 二次不等式：看拋物線}

解 $ax^2+bx+c>0$（或 $<0$）的標準流程：
\begin{enumerate}\setlength{\itemsep}{1pt}
\item 先讓領導係數為正（若 $a<0$，兩邊乘 $-1$ 並把不等號反向），畫出\textbf{開口向上}的拋物線。
\item 因式分解或用公式求兩根 $\alpha\le\beta$（拋物線與 $x$ 軸的交點）。
\item 開口向上時，圖形在兩根\textbf{外側}為正、\textbf{中間}為負。於是 $>0$ 取兩側（$x<\alpha$ 或 $x>\beta$）、$<0$ 取中間（$\alpha<x<\beta$）。
\end{enumerate}

\begin{補充框}{「大於取兩邊、小於取中間」從何而來？}
這句口訣是\textbf{圖形的結論}。令 $f(x)=ax^2+bx+c$，解不等式本質上是在問「拋物線什麼時候在 $x$ 軸上方／下方」。而拋物線\textbf{只在根處穿過 $x$ 軸、由一側換到另一側}，所以兩個相異根 $\alpha<\beta$ 把 $x$ 軸切成三段，每段拋物線\textbf{整段都在同一側}。剩下的只看開口方向：
\begin{itemize}\setlength{\itemsep}{1pt}
\item \textbf{開口向上（$a>0$）}：兩端翹起、中間凹下，故\textbf{兩根外側在上方（$>0$）、中間在下方（$<0$）}，即「大於取兩邊、小於取中間」。
\item \textbf{開口向下（$a<0$）}：上下顛倒，口訣會反過來。這也是課本要求\textbf{先把開口調成向上}的原因，以便統一使用同一句口訣。
\end{itemize}
也可純代數驗證：開口向上時 $f(x)=a(x-\alpha)(x-\beta)$，$a>0$。$x$ 在兩根之間時 $(x-\alpha)>0$、$(x-\beta)<0$，乘積為負，$f(x)<0$；$x$ 在兩根之外時兩因式同號，乘積為正，$f(x)>0$。純算結果與看圖一致。
\end{補充框}

上面「大於取兩邊、小於取中間」只處理了「有兩個相異根」的情形。但兩根會不會重合、甚至根本沒有實根，取決於\textbf{判別式（discriminant）$D=b^2-4ac$}。解二次不等式前，務必先看拋物線\textbf{碰不碰得到 $x$ 軸}，因為這會把解的形態整個改變。下面統一在\textbf{開口向上（$a>0$）}下分三種情況（若 $a<0$，先乘 $-1$ 並把不等號反向，調成向上）。

\begin{補充框}{依判別式 $D$ 分三種情況（開口向上 $a>0$）}
\textbf{一、$D>0$（兩相異根 $\alpha<\beta$）}：拋物線穿過 $x$ 軸兩次。$>0$ 取兩根外側（$x<\alpha$ 或 $x>\beta$）、$<0$ 取中間（$\alpha<x<\beta$），即前述口訣。

\textbf{二、$D=0$（重根 $\alpha$）}：拋物線與 $x$ 軸\textbf{相切}於 $x=\alpha$（碰一下而不穿過），此時 $f(x)=a(x-\alpha)^2\ge 0$ 恆成立。四種不等號的解都不是區間：
\begin{itemize}\setlength{\itemsep}{1pt}
\item $f(x)>0$：除切點外恆成立 $\Rightarrow$ 解為 $x\neq\alpha$（即除 $\alpha$ 以外的所有實數）。
\item $f(x)\ge 0$：\textbf{所有實數}（含切點）。
\item $f(x)<0$：\textbf{無解}（拋物線從不在 $x$ 軸下方）。
\item $f(x)\le 0$：\textbf{只有 $x=\alpha$ 一點}（只有切點讓 $f=0$）。
\end{itemize}

\textbf{三、$D<0$（無實根）}：拋物線整條在 $x$ 軸上方，恆 $>0$。$f(x)>0$（或 $\ge 0$）的解是\textbf{所有實數}；$f(x)<0$（或 $\le 0$）\textbf{無解}。

\textbf{注意：}遇到 $a<0$ 不可直接套「大於取兩邊」，須先確認開口方向或統一乘 $-1$ 調成向上。$D=0$ 與 $D<0$ 時答案不是「某個區間」，而是「全體實數」「無解」或「單獨一點」；看到 $\le,\ge$ 還要留意切點那一點要不要算進去。
\end{補充框}

\fig{d12-quad-ineq}{0.92}{圖 2-4：以 $x^2-x-6=(x-3)(x+2)$（根 $-2,3$，開口向上）為例。左——$x^2-x-6>0$ 取圖形在 $x$ 軸上方者，解為 $x<-2$ 或 $x>3$（兩邊）；右——$x^2-x-6<0$ 取下方者，解為 $-2<x<3$（中間）。}

\subsection*{\sffamily C.\ 已分解的高次不等式：符號分析}

若不等式已分解成一串一次因式相乘，例如 $(x-1)(x+2)(x-3)>0$，就用\textbf{數線分段定號}：把根 $-2,1,3$ 標在數線上切成四段，在每段任取一個測試點代入，看整體乘積的正負號，再取符合方向的區段。

當根都是\textbf{單根}時，正負號\textbf{每過一個根就翻一次}，因此可從最右段（必為 $+$）往左交替判斷，不必逐段計算。

\textbf{注意：}
\begin{itemize}\setlength{\itemsep}{1pt}
\item 不等式\textbf{不能像方程式一樣交叉相乘消去未知數}（不知未知量正負，乘過去不知是否該變號）。一律移項成「$\cdots>0$」再分析。
\item 端點是否包含看不等號：嚴格 $>,<$ 不含端點（空心），$\ge,\le$ 含端點（實心）。但若某點使分母為零或式子無意義，無論如何都要排除。
\item \textbf{重根會「不變號」}：如 $(x-1)^2(x-3)>0$，過 $x=1$ 時號不翻（平方恆非負），須特別處理。
\end{itemize}

\section*{\sffamily 2-5\quad 本章重點整理}

\begin{補充框}{一頁複習（公式與關鍵結論）}
\begin{itemize}\setlength{\itemsep}{2pt}
\item \textbf{除法原理}：$f(x)=g(x)q(x)+r(x)$，$\deg r<\deg g$。
\item \textbf{餘式定理}：$f(x)$ 除以 $x-a$ 的餘式 $=f(a)$。\textbf{因式定理}：$x-a$ 為因式 $\iff f(a)=0$。兩者皆由除法原理代 $x=a$ 推得。
\item \textbf{綜合除法}：除以 $x-a$ 時「落、乘、加」搬係數（缺項補 $0$），最右一格為 $f(a)$。
\item \textbf{一次函數}：$f(x)=mx+k$，斜率 $m$；與 $y=mx$ 只差鉛直平移 $k$。\textbf{分點公式}：$P=\dfrac{na+mb}{m+n}$。
\item \textbf{配方法}：$ax^2+bx+c=a(x-h)^2+k$，$h=-\dfrac{b}{2a}$，$k=f(h)$。頂點 $(h,k)$、對稱軸 $x=h$；$a>0$ 最小值 $k$、$a<0$ 最大值 $k$。圖形由 $y=ax^2$ 平移到頂點。
\item \textbf{閉區間最值}：在 $[p,q]$ 上先看對稱軸 $h$ 在不在區間內；在內則一端是頂點、另一極值在離 $h$ 較遠的端點；不在則最值都在兩端點 $f(p),f(q)$。
\item \textbf{三次函數}：對稱中心 $x=-\dfrac{b}{3a}$（點對稱）；大域看最高次項，局部近似直線；至少有一個實根。
\item \textbf{多項式不等式}：找根、切數線、定號取段。二次（開口向上）依判別式分情況：$D>0$「大於取兩邊、小於取中間」；$D=0$ 相切，解為「除切點外」「全體」「單點」或「無解」；$D<0$ 不碰 $x$ 軸，恆正 $\Rightarrow$ 全體實數或無解。高次：單根變號、重根不變號。
\end{itemize}
\end{補充框}

\end{document}
